\chapter{Application Generation Vision}
\label{chap:application-vision}

The server-class cluster vision (\autoref{chap:cluster-vision}) asks what an operating system
might look like if signed execution objects, explicit capabilities, and
cluster-level placement were native mechanisms rather than layers constructed
above general-purpose Linux. That direction creates an immediate application
problem: a deliberately small, non-POSIX userspace does not run the existing
body of server software, and asking ordinary application developers to program
directly against a new kernel interface would make the platform impractical.

This chapter develops a proposed answer. A process-oriented generator such as
Durga (\url{https://github.com/FrodeRanders/durga}) can mediate between
application intent and the Charlotte execution environment. Developers describe
a process in BPMN and implement business rules against a narrow generated
contract. A Charlotte backend supplies executable wrappers, capability
requirements, communication structure, and cluster desired state. The tools and
domain model absorb most platform-specific knowledge.

The implementation already crosses both project boundaries. Durga's
\texttt{charlotte} target emits a \texttt{no\_std} adapter, Cargo and AArch64
build inputs, capability requests, named Kafka connector profiles, explicit
stack and thread limits, signed replica and affinity policy, and deployment
metadata for each BPMN component. A developer-owned
\texttt{charlotte/resources.yaml} is created once, preserved across later
regeneration, and must be reviewed before deployment output becomes usable. The
same retained values drive descriptor-signing options.

Generated adapters compile into Charlotte ELFs. After CI builds, signs, and
uploads them, Charlotte admits their descriptors atomically in a signed
\texttt{CRELEASE}, resolves replica placement through Raft, fetches artifacts
from central S3, mediates grants, launches isolated domains, and reports
exact-generation readiness. Business handlers and schemas remain
developer-owned completion points. A richer semantic bundle that binds the BPMN
model, schemas, and provenance, together with a production release controller,
remains the next boundary.

\architecturewidefigure{durga-charlotte-generation}
    {Durga turns process intent into Charlotte-native adapters, capability and
    connector requirements, signed release inputs, and cluster desired state.}
    {fig:durga-charlotte-generation}

\section{The application thesis}

The conventional response to a new operating system is to reproduce a familiar
application interface. POSIX compatibility, a C library, sockets, a filesystem
hierarchy, dynamic linking, language runtimes, and ports of common packages
provide a path for existing programs. This effective compatibility also
recreates execution assumptions from which Charlotte is trying
to depart: ambient access to a machine-like namespace, software assembled from
mutable packages, and a process API designed around one host.

The alternative considered here is compatibility at a higher level:

\begin{center}
\begin{tikzpicture}[
    box/.style={draw, rounded corners, align=center, inner sep=5pt},
]
    \node[box] (b) {business process\\+ business rules};
    \node[box, below=0.6cm of b] (g) {generator};
    \node[box, below=0.6cm of g] (n) {native Charlotte\\application};
    \draw[->] (b) -- (g);
    \draw[->] (g) -- (n);
\end{tikzpicture}
\end{center}

Compatibility moves into libraries, data formats, external-system connectors,
diagnostics, testing, and developer tools. For a useful class of
process-oriented business applications, developers may never need to see the
operating-system interface. The
generator and its runtime support can be the only code that understands such
details as message endpoints, capability handles, service registration, memory
objects, lifecycle events, supervision, and deployment metadata.

The resulting thesis complements \autoref{chap:cluster-vision}:

\begin{quote}
A constrained execution environment need not imply a difficult application
environment when infrastructure code is generated from a higher-level model.
CharlotteOS need not expose a familiar operating-system API if familiar tools
can generate Charlotte-native applications.
\end{quote}

The scope is deliberately narrower than all applications or existing software:
it covers domains where the process graph is already an important source of
truth and the surrounding infrastructure is regular enough to generate.

\section{Durga's present application model}

Durga is presently a BPMN-driven Kafka scaffolding and monitoring system. The
BPMN model is executable input to a generator rather than documentation added
after implementation. Durga derives workers, gateways, orchestration, topic
setup, process lifecycle events, and monitoring structure from that model. The
application developer completes or supplies the task-specific business
behavior while generated code implements much of the surrounding protocol.

Three current backends demonstrate that the process description need not be tied
to one implementation language:

\begin{itemize}
    \item The Java target generates a Quarkus application using SmallRye
        Reactive Messaging. Tasks, gateways, and orchestration become CDI
        components wired to Kafka channels, with shared runtime support.

    \item The Rust target generates a Cargo crate. Tasks and gateways become
        binaries under \texttt{src/bin}, with common code in a generated
        library and shared Durga runtime support.

    \item The prototype Charlotte target generates minimal \texttt{no\_std}
        activity adapters and a deployment plan. Kafka broker credentials stay
        in separately named connector profiles; applications receive only
        granted service endpoints.
\end{itemize}

The Java and Rust targets use the same process-event wire model and topic
layout, so their components interoperate through the generated protocol. The
Charlotte target consumes the same process topology and generates Kafka
connector intent. Its current adapters establish the platform and deployment
boundary; the complete process-event runtime is the next semantic layer. The
process model can therefore remain stable while code generation changes the
language and runtime used to realize it.

The Java and Rust deployment paths remain conventional. A generated application
can include a JAR or Rust binaries, Kafka configuration and topic provisioning,
a container build, a Kubernetes Deployment and Service, health probes, resource
limits, environment configuration, and deployment scripts. Scaling combines
application replicas with Kafka consumer groups and topic partitions. In an
OpenShift environment, image policy, registry access, security context,
network policy, routes, secrets, monitoring integration, and organizational
admission policy normally add further deployment concerns. The Charlotte
target instead emits platform build and deployment-planning inputs; operating
the release remains outside Durga.

These mechanisms supply real isolation, durability, routing, scaling, recovery,
and operational control for conventional applications on general-purpose
hosts. Their distance from the original business process motivates a more
direct generated representation for Charlotte.

\section{Three distinct languages}

The proposed architecture is easier to reason about if three layers are kept
separate.

\subsection{The business and process language}

The upper layer contains BPMN and developer-owned business rules. It describes
activities, gateways, events, participants, message flows, data objects, timers,
failures, and process completion. Handwritten code supplies domain behavior
such as validating an invoice, calculating an entitlement, or selecting a
payment route.

This layer should remain largely independent of Charlotte. A developer should
understand the business domain, the supported BPMN subset, and the Durga task
contract. Charlotte-specific message buffers, deployment calls, signing notes,
and capability handles should not leak into ordinary rule code unless the
application deliberately requests a platform-specific facility.

\subsection{The Durga semantic execution model}

The middle layer gives execution meaning to the process. Its concepts include:

\begin{itemize}
    \item process definitions and versioned process instances;
    \item activity entry, completion, failure, cancellation, and retry;
    \item gateways, branch selection, fan-out, joining, and correlation;
    \item messages, signals, timers, call activities, and subprocess scopes;
    \item typed or schema-identified inputs and outputs;
    \item state transitions and lifecycle observations; and
    \item delivery, replay, idempotency, and failure semantics.
\end{itemize}

Some of these semantics are explicit in BPMN. Others are Durga runtime choices
or policies that must be made explicit before an alternative backend can claim
equivalence. For example, a sequence flow orders two activities. Delivery
durability, repeated execution, acknowledgement, and recovery of a partially
completed external side effect require additional explicit policy.

The semantic model must therefore be a real intermediate representation rather
than a renaming of Kafka topics. A Charlotte backend should consume concepts
such as activity, event, correlation, retry, and durable boundary.

\subsection{The execution substrate}

The lower layer realizes the semantic model. The established Durga backends use
Java or Rust workers and Kafka communication. The prototype Charlotte target
starts to use Charlotte execution objects, capabilities, names, and cluster
desired state, while retaining an external event log where its semantics are
required. Supervision and the complete process runtime remain target work.

The intended relationship is therefore:

\begin{center}
\begin{tikzpicture}[
    box/.style={draw, rounded corners, align=center, inner sep=5pt},
]
    \node[box] (root) {BPMN + business rules};
    \node[box, below=0.7cm of root] (mid) {Durga semantic\\process model};

    \node[box, below left=1.2cm and 1.4cm of mid] (java) {Java/Kafka};
    \node[box, below=1.2cm of mid] (rust) {Rust/Kafka};
    \node[box, below right=1.2cm and 1.4cm of mid] (char) {Charlotte\\backend};

    \draw[->] (root) -- (mid);
    \coordinate (m) at ([yshift=-0.5cm]mid.south);
    \draw (mid.south) -- (m);
    \draw (java.north |- m) -- (char.north |- m);
    \draw (java.north) -- (java.north |- m);
    \draw (rust.north) -- (rust.north |- m);
    \draw (char.north) -- (char.north |- m);
\end{tikzpicture}
\end{center}

Each lower-layer mapping must preserve the process contract. When the available
substrate lacks a required semantic, the generator retains a suitable service
such as Kafka or rejects the build. Moving bytes between activities is only one
part of backend correctness.

\section{The Charlotte backend and its target architecture}

The working generator establishes the source, build, authority, connector, and
deployment skeleton described at the start of this chapter. The target
architecture fills that skeleton with completed business-rule bindings and a
semantic process runtime, then adds bundle-level admission, resource-aware
rolling placement, lifecycle supervision, and rollback. Charlotte's existing
release path already supplies fixed replica sets and exact-generation DSR
readiness below that semantic layer.

\subsection{Generated activity boundary}

For a simple activity, the developer-owned interface might be no broader than:

\begin{lstlisting}[style=rustCode]
fn calculate_entitlement(
    claim: Claim
) -> Result<Decision, RuleError>
\end{lstlisting}

The exact Rust interface is a design question, not an implemented ABI. Around
that function, the Charlotte backend could generate a wrapper that:

\begin{enumerate}
    \item receives a typed activity invocation through a delegated endpoint;
    \item validates the process, activity, instance, schema, and correlation
        identities carried by the invocation;
    \item deserializes input into the developer-facing type;
    \item emits an activity-entered observation;
    \item invokes the business rule within its assigned protection domain;
    \item converts success or failure into the Durga execution result;
    \item serializes the result and sends it only through generated output
        capabilities;
    \item emits completion, failure, and resource observations; and
    \item participates in generated supervision, shutdown, and upgrade
        protocols.
\end{enumerate}

The generated wrapper owns platform mechanics: receiving Charlotte messages,
accepting capabilities, allocating shared memory, re-registering after
replacement, responding to endpoint closure, and attaching telemetry. The
developer-facing business contract remains suitable for unit tests on an
ordinary development host.

The source layout enforces the same boundary. Regeneration replaces generated
files and preserves developer-owned business modules. Specialized services may
use explicit escape hatches; ordinary rules remain free of raw Charlotte
handles.

\subsection{A deliberately narrow runtime interface}

The generated runtime would target a small Charlotte service interface rather
than POSIX. A candidate operation set might include:

\begin{itemize}
    \item receive from and send to delegated message endpoints;
    \item allocate, transfer, and release bounded memory objects;
    \item access a delegated clock, random source, or cryptographic service;
    \item invoke a named service through an explicit capability;
    \item emit structured observations;
    \item checkpoint explicitly owned state; and
    \item cooperate with supervision, replacement, and cancellation.
\end{itemize}

This list is proposed. The experiment should discover the minimum viable
interface rather than declare it in advance. In particular, external protocols,
TLS, databases, object storage, and durable logs will require trusted connectors
or Charlotte-native services. Code generation cannot make those integrations
disappear; it can make authority to use them explicit and can keep their
mechanics outside the business rule.

\section{From process graph to capability graph}

Charlotte workloads begin with no ambient authority and receive only delegated
capabilities. Durga has information from which a first capability request can be
generated: which activity receives which input, which activity can send to
which successor, which participant communicates with which other participant,
and which declared data or external service an activity uses.

For a simple process fragment:

\begin{center}
\begin{tikzpicture}[
    act/.style={draw, rounded corners, align=center, inner sep=5pt, font=\small},
]
    \node[act] (vc) {ValidateClaim};
    \node[act, right=0.8cm of vc] (ce) {CalculateEntitlement};
    \node[act, right=0.8cm of ce] (id) {IssueDecision};
    \node[act, below=0.8cm of ce] (rs) {RulesService};

    \draw[->] (vc) -- (ce);
    \draw[->] (ce) -- (id);
    \draw[->] (ce) -- (rs);
\end{tikzpicture}
\end{center}

the backend could derive a graph in which:

\begin{itemize}
    \item \texttt{ValidateClaim} may send validated claims to
        \texttt{CalculateEntitlement};
    \item \texttt{CalculateEntitlement} may invoke
        \texttt{RulesService} and send decisions to
        \texttt{IssueDecision};
    \item \texttt{IssueDecision} receives decisions but cannot invoke the
        rules service; and
    \item no activity receives a generic network, filesystem, or cluster
        administration capability.
\end{itemize}

The process communication graph can thus become an input to the Charlotte
capability graph. BPMN message flows, pools, participants, and explicitly
annotated external interactions are particularly useful because they already
express architectural boundaries. A message between a claims pool and a
payments pool may justify an explicit cross-domain endpoint rather than ambient
reachability between their executables.

Absence of a BPMN edge is useful negative information. A complete authority
decision also accounts for implicit gateway and orchestration paths, monitoring
and administration channels, and external systems declared by business-rule
code. Capability inference therefore combines several sources:

\begin{itemize}
    \item the supported BPMN topology and its scopes;
    \item typed task contracts and data declarations;
    \item explicit annotations for external services and durable resources;
    \item generator-known runtime and observation requirements;
    \item optional static analysis of the developer-owned code; and
    \item organization policy that may narrow, approve, or reject the inferred
        request.
\end{itemize}

The generator produces a reviewable authority request. A separate build or
admission policy decides which capability bindings are acceptable. This retains the distinction from
\autoref{chap:cluster-vision} between
an artifact's signed claims and the cluster authority that admits and delegates
them.

\subsection{Security graph and communication graph}

The execution graph, communication graph, and security graph are related but
not identical:

\begin{itemize}
    \item The execution graph says which semantic transition may follow which
        other transition.

    \item The communication graph says which generated component must exchange
        which message or object with which other component.

    \item The security graph says which protection domain is authorized to
        perform that exchange at runtime.
\end{itemize}

One logical sequence flow need not create one network-visible endpoint, and two
activities fused into the same executable may still be kept behind separate
logical interfaces for observation or future decomposition. Conversely, a
shared orchestrator may require communication not drawn as a direct business
edge. Keeping the three graphs distinct prevents a convenient compilation
choice from silently becoming the security policy.

\section{The signed process bundle}

The proposed backend would produce more than executable code. Its natural
deployment unit is a signed process bundle that preserves the relationship
between model, rules, generated wrappers, and cluster policy. A bundle could
contain:

\begin{itemize}
    \item the original BPMN model and a canonical digest of the model;
    \item the Durga semantic intermediate representation and its version;
    \item one or more Charlotte-native executable artifacts;
    \item immutable identities and signatures for every executable;
    \item input, output, event, and state schemas with stable identities;
    \item the execution, communication, and requested capability graphs;
    \item signer, build, generator, source, dependency, and toolchain
        provenance;
    \item an SBOM or the signed digest of retained SBOM and attestation
        evidence;
    \item resource estimates, concurrency declarations, replication policy,
        placement hints, and failure-domain constraints;
    \item durable-log and external-connector requirements;
    \item observation metadata mapping runtime identities back to BPMN
        activities; and
    \item the proposed cluster desired state for this process release.
\end{itemize}

Execution resources and replica intent deserve special treatment among those estimates. Only the application
and generator authors know the component's worst-case call depth, generated
adapter frames, stack-resident data, safe concurrency, and intended replication. Durga therefore
creates a developer-owned \texttt{charlotte/resources.yaml} carrying explicit
\texttt{stackPagesPerThread}, \texttt{maxThreads}, and
\texttt{shutdownGraceMillis} requests together with replica count,
every-eligible-node selection, spread, affinity, and anti-affinity. It creates that
file once and never overwrites it during regeneration. Operations may approve
or reject the reviewed values when signing the deployment, but should not have
to guess them from the ELF or silently rewrite them. Charlotte's
\texttt{CDEPLOY5} descriptor signs all three values and placement policy; Durga
derives the corresponding signing options and adds the CLS2 parallel-instance
requirement whenever the policy permits concurrent instances. The kernel gives every thread
in the protected domain the specified 4 KiB-page stack allocation and counts
the bootstrap thread against the active-thread quota. It first requests
cooperative application draining and then forcibly retires the domain at the
signed deadline. Both execution-resource ranges are one through 64; shutdown
grace is zero through 300,000 milliseconds. A quota-exceeding spawn aborts the domain under the
current fail-closed ABI. Legacy \texttt{CDEPLOY1} descriptors retain the former
four-page and 16-thread defaults; \texttt{CDEPLOY2} retains its signed stack
value and receives the 16-thread compatibility default; and \texttt{CDEPLOY3}
retains both execution limits and receives the five-second shutdown default.

\autoref{chap:cluster-vision} already implements signed, name-bound ELF artifacts with release,
rollback, class, flags, and an optional provenance digest, and it pins the
complete ELF digest in a deployment record. A process bundle would build on
those mechanisms; it would not replace executable verification with trust in a
container or archive. Each executable remains independently verifiable, while
the signed bundle states that a specific set of artifact and schema identities
collectively realizes one version of a process.

Bundle signing is proposed future work. It also introduces supply-chain
questions. The generator and build toolchain become security-sensitive because
they derive authority and create wrappers. Reproducible generation, separation
between requested and granted capabilities, independent policy review, retained
provenance, and verification of the bundle-to-artifact links are necessary.
Signing generated output does not prove that the BPMN model or generator was
correct.

\subsection{Generated desired state}

The bundle could carry a declarative target such as:

\begin{verbatim}
process: claims
release: 42

component validate:
    artifact: sha256:...
    stack-pages-per-thread: 4
    max-threads: 4
    shutdown-grace-millis: 15000
    replicas: 4
    parallel: true

component calculate:
    artifact: sha256:...
    stack-pages-per-thread: 8
    max-threads: 8
    shutdown-grace-millis: 30000
    replicas: 8
    parallel: true
    requires: service:entitlement-rules

placement:
    affinity: [validate, calculate]
    minimum-failure-domains: 2

durability:
    ClaimSubmitted: external-event-log
    validate-to-calculate: native-message
\end{verbatim}

The example illustrates deployment intent generated from the same semantic
model as the executable wrappers; a future format may use different syntax.
Cluster admission can validate the requested replica counts,
parallel-safety declaration, artifact identities, stack limits, capabilities,
and placement constraints before committing desired state to the replicated
manifest.

The Charlotte manifest accepts signed singleton, fixed-replica, and
\emph{every eligible node} policy. The \texttt{PlacementPolicy} contract
described in \autoref{chap:cluster-vision} is stored in \texttt{CDEPLOY5}; a
leader-side controller resolves it and commits concrete replica sets. A
generated component set can be wrapped in \texttt{CRELEASE}, atomically
admitted, and observed together with \texttt{cluster-sign release-apply}.
This binds exact executable deployment descriptors and placement policy. The
semantic bundle adds the BPMN digest, schemas, communication graph, and
provenance. Consuming it also depends on resource-aware scheduling and the
production administration authority identified there.

\section{Process topology as placement information}

A general-purpose cluster scheduler normally begins with deployable units,
resource requests, and manually supplied affinity rules, then infers
relationships from traffic and resource telemetry. Durga can additionally
supply the application's semantic process graph before execution.

Durga knows useful structure before execution begins:

\begin{itemize}
    \item which activities are adjacent and expected to exchange data;
    \item where the graph fans out, joins, waits, or selects one branch;
    \item which activities may execute concurrently;
    \item which events enter from or leave for external systems;
    \item which subprocesses and participants form natural boundaries; and
    \item which activities are expected to block on timers, people, or durable
        external events.
\end{itemize}

A Charlotte backend can translate this into an initial placement hypothesis.
Frequently adjacent activities may receive a common affinity group; replicas of
one parallel-safe activity may receive anti-affinity or minimum-failure-domain
rules; an activity using a restricted service may be constrained to the same
trust domain; a high-volume data edge may favor co-location.

The BPMN topology supplies a hypothesis rather than a measurement. A rarely
selected branch may carry very large objects, an adjacent pair may spend almost
all its time in an external service, and one process definition may serve
different data distributions. Runtime evidence must refine incomplete or
incorrect resource estimates and expected frequencies.

Charlotte's observation substrate can refine it. Runtime identities generated
from the process model allow the cluster to measure message counts, bytes,
latency, queue pressure, CPU time, memory pressure, failures, and retry rates per
activity and edge. A future placement controller can compare declared structure
with observed behavior, then move or replicate components when sustained
evidence justifies the change. Availability and trust constraints remain hard
constraints; observed affinity must not collapse replicas into one failure
domain.

This gives the process compiler and cluster different roles:

\begin{verbatim}
Durga model       -> initial graph and declared constraints
Charlotte runtime -> measured graph and resource evidence
placement policy  -> admissible movement and replication
controller        -> assignments in replicated cluster state
\end{verbatim}

Charlotte already resolves declared replica and affinity policy into automatic
assignments. The proposed feedback loop adds cluster-wide observation and uses
sustained runtime evidence to refine those assignments.

\section{Logical activities and physical executables}

BPMN describes logical units of work independently of protection domains,
operating-system processes, containers, or scheduled artifacts. The current
Rust generator maps activities and gateways to binaries. A Charlotte backend
may instead choose physical granularity from authority, durability, placement,
and performance constraints.

For a logical chain:

\begin{verbatim}
A -> B -> C
\end{verbatim}

the backend might generate any of:

\begin{verbatim}
[A] -> [B] -> [C]
[A B] -> [C]
[A B C]
\end{verbatim}

Fusion can remove serialization, endpoint, scheduling, and network overhead. It
may also permit intermediate values to remain in one address space. Separation
supports independent scaling, fault containment, least authority, different
trust domains, independent upgrades, and placement near distinct resources.

Executable granularity should therefore be an explicit compiler and policy
decision. Relevant constraints include:

\begin{itemize}
    \item capability and participant boundaries;
    \item durable handoff points and replay requirements;
    \item independent scaling and parallel-safety declarations;
    \item failure and retry domains;
    \item state ownership and upgrade boundaries;
    \item expected message size and frequency;
    \item architecture or hardware requirements; and
    \item the cost of losing a fused group when one activity fails.
\end{itemize}

The logical activity identity must survive fusion. Observations, failures,
schema versions, and business audit events should still be attributable to the
corresponding BPMN activity. The signed bundle should state the mapping from
logical activities to physical artifacts, so that the cluster can explain its
placement and so that a later release can split or fuse components without
renaming the business process.

Automatic re-fusion from runtime evidence is a later research topic. Because
executable boundaries affect security, state, and upgrade semantics, the first
implementation should require a new generated and signed release for every
change.

\section{Native messages and durable event logs}

Kafka presently supplies more than byte transport to Durga. It offers durable
logs, replay, partition ordering, consumer groups, offset tracking, integration
with external producers and consumers, and transactional options. A Charlotte
message endpoint must not be treated as equivalent merely because it can carry
the same payload.

CharlotteOS now has an implemented, capability-scoped Kafka compatibility
service (\autoref{sec:kafka-service}). A provisioned instance exposes an
explicit broker allow-list, a fixed consume partition, up to 64 allow-listed produce routes, a
group, transactional identity, and separately named authority endpoints. Its
owned Rust client supports idempotent produce, bounded read-committed consume,
explicit delivery acknowledgement, and a transaction that atomically combines
Kafka output with the consumed delivery's group-offset advance. This supplies
a concrete external event-log boundary for a Durga backend.

The generated deployment should treat the broker connector and the behavioural
role as separate authorities. Only \texttt{kafka} receives the immutable,
digest-checked profile containing endpoint, trust, and authentication
secrets. A producer, consumer, or transactional-step runner receives a
connection capability to the corresponding profile-defined access point. The
connector enforces that endpoint's Kafka-rights ceiling for every opcode, so a
producer-only connection cannot be used to consume or transact. The activity procedure behind a step
receives only the bounded invocation capability and no Kafka authority.
Each connector profile fixes all bounded access-point names, so a deployment
can provision several least-authority connections from one secret-bearing
connector without replacing a global \texttt{kafka} publication. The runner
resolves only the transactional access point declared by its step profile.
The generic \texttt{kafka\_step} runner now implements this split: it calls the
procedure with borrowed input and an attempt number, validates its bounded
output batch, and atomically produces admitted outputs with the consumed
offset. It retries timeouts and explicit transient failures and writes terminal,
malformed, or exhausted inputs to a configured DLQ. The remaining deployment
work is direct capability injection, generated artifact construction, and
controller-managed lifecycle/fencing rather than the transaction loop itself.

The useful distinction is between application events whose meaning is
intrinsically durable and internal execution plumbing whose purpose is to
advance a running process.

An externally meaningful event such as \texttt{ClaimSubmitted} may need a
retained, replayable, independently consumable record. Kafka or another durable
event-log service may remain the correct substrate. An internal instruction
such as ``activity A completed; make B runnable'' may need only the delivery,
deduplication, and recovery guarantees of the Durga runtime and could use native
Charlotte messaging.

A proposed Charlotte deployment may therefore be hybrid:

\begin{center}
\begin{tikzpicture}[
    box/.style={draw, rounded corners, align=center, inner sep=5pt},
    act/.style={draw, rounded corners, align=center, inner sep=4pt},
]
    \node[box] (ext) {external producers};
    \node[box, below=0.7cm of ext] (log) {durable event log};

    \node[act, below=1.2cm of log] (c) {C};
    \node[act, left=0.5cm of c] (b) {B};
    \node[act, left=0.5cm of b] (a) {A};
    \node[align=center, right=0.25cm of c] (note) {Native\\Charlotte\\messages};

    \node[draw, left=0.0cm of a, dashed, rounded corners, fit=(a)(b)(c)(note), inner sep=8pt] (group) {};

    \node[box, below=1.2cm of c] (bev) {durable business event};

    \draw[->] (ext) -- (log);
    \draw[->] (log) -- (group);
    \draw[->] (a) -- (b);
    \draw[->] (b) -- (c);
    \draw[->] (group) -- (bev);
\end{tikzpicture}\end{center}

This keeps Kafka where log semantics and ecosystem integration are required,
while avoiding a broker round trip for every internal transition. It also makes
the durable boundary visible in the process bundle and capability request.

The present Kafka adapter is intentionally narrower than a general client. It
uses a profile-fixed partition and a Charlotte fixed-partition group assignor
while routing among an explicit broker allow-list. It supports group join,
heartbeat, standby failover, and generation-fenced ordinary and transactional
offset commits, but not conventional/cooperative assignor interoperability,
complete coordinator-migration recovery, client-certificate chains, or
compression. It supports verified TLS,
connector-only SCRAM-SHA-256 authentication, and connector-only P-256 mTLS,
independently or together. A generated deployment
would therefore need either to accept that bounded profile, extend the adapter
under equally explicit authority, or retain a conventional connector outside
Charlotte for requirements the native service cannot satisfy.

The distinction cannot be inferred reliably from sequence flow alone. The
process author or policy must declare durability, ordering, replay, retention,
and delivery requirements. The Charlotte runtime must define how it recovers an
in-flight process when a node fails, when an acknowledgement is lost, or when a
business rule completes twice. If native messaging does not yet satisfy those
requirements, the backend must retain Kafka for that edge. Removing Kafka from
internal paths is an optimization and architectural option, not a prerequisite
for the Charlotte backend.

\section{Preserving application intent}

In a conventional deployment, semantic information is progressively lowered
into infrastructure objects. A BPMN activity may eventually appear to the
cluster as a container image, a Deployment, CPU and memory values, labels, and
network endpoints. Those objects are operationally useful, but they retain
little explanation of why two components communicate or what a latency
measurement means to the business process.

Durga and Charlotte together could preserve an identity chain:

\begin{center}
\begin{tikzpicture}[
    box/.style={draw, rounded corners, align=center, minimum width=6cm, inner sep=5pt},
]
    \node[box] (id1) {BPMN process and\\activity identity};
    \node[box, below=0.5cm of id1] (id2) {Durga execution and\\schema identity};
    \node[box, below=0.5cm of id2] (id3) {capability and\\communication edges};
    \node[box, below=0.5cm of id3] (id4) {physical artifact and\\placement identity};
    \node[box, below=0.5cm of id4] (id5) {runtime observations and\\process outcomes};

    \draw[->] (id1) -- (id2);
    \draw[->] (id2) -- (id3);
    \draw[->] (id3) -- (id4);
    \draw[->] (id4) -- (id5);
\end{tikzpicture}
\end{center}

Every generated message and observation can carry or imply stable process,
release, instance, activity, edge, and schema identities. The signed bundle
binds those logical identities to executable digests. The cluster desired state
binds artifact identities to admitted capability and placement policy. Runtime
telemetry can then be mapped back to the BPMN diagram.

This allows questions at several levels to share one evidence chain:

\begin{itemize}
    \item Which business activity accounts for the current queue pressure?
    \item Which capability edge carried the slow calls?
    \item Which artifact release implemented that activity?
    \item On which nodes and trust domains were its instances placed?
    \item Did the observed traffic support the generated affinity hint?
    \item Did an executable fusion decision improve latency while preserving
        the intended fault and security boundaries?
\end{itemize}

Durga already emits lifecycle events and uses the BPMN model as the background
for process monitoring. Charlotte already exposes executor observations and
generation-aware names. Joining those identities and feeding the result into
cluster placement is proposed future work.

Telemetry must not become ambient surveillance authority. Generated observation
channels should be explicit capabilities; payloads should avoid business data
unless required and authorized; and retention and access belong to policy.
Preserving intent is valuable only if the observation path preserves the same
security discipline as execution.

\section{Comparison with Kubernetes and OpenShift}

The comparison asks which layer should own each orchestration responsibility
and which deployment objects the generator should produce.

In the current conventional path, Durga may generate or depend on:

\begin{verbatim}
BPMN + rules
    -> Java/Rust application
    -> Kafka topics and configuration
    -> JAR or native executable
    -> container filesystem image
    -> image registry and image policy
    -> Kubernetes/OpenShift manifests
    -> Pods, Services, probes and NetworkPolicy
    -> container runtime
    -> Linux isolation and resource mechanisms
\end{verbatim}

In the proposed Charlotte path:

\begin{verbatim}
BPMN + rules
    -> Durga semantic model
    -> signed process bundle
    -> Charlotte desired state
    -> Charlotte cluster
\end{verbatim}

The shorter Charlotte path assigns each concern to the operating system, the
signed bundle, or an explicit external service:

\begin{itemize}
    \item An OCI filesystem image may become a signed, self-contained Charlotte
        execution artifact. Artifact storage, distribution, verification, and
        rollback remain necessary.

    \item A Kubernetes Deployment may become Charlotte replicated desired state
        plus replica and placement policy. Reconciliation, failure recovery,
        rollout, and admission remain necessary.

    \item A Service and service account may become a Charlotte name and
        explicitly delegated invocation capability. Naming, routing, identity,
        and key management remain necessary.

    \item NetworkPolicy may be replaced for generated internal communication by
        the admitted capability graph. Gateways to IP networks and external
        systems still require network security policy.

    \item Liveness and readiness endpoints may become native supervision and
        generated lifecycle state. The cluster must still distinguish work
        that is slow, blocked, failed, or safe to restart.

    \item Pod affinity may begin as a generated process communication graph and
        later be refined by Charlotte observation. Scheduling and failure-domain
        policy remain necessary.

    \item Charlotte may not need a container runtime or container-specific
        Linux policy. Namespaces, cgroups, and seccomp may also be absent because
        Charlotte starts with separate address spaces and delegated authority.
        Equivalent memory, CPU, and device accounting must nevertheless exist
        natively.

    \item Kafka topics used only for internal execution plumbing may become
        Charlotte message endpoints. Kafka remains where durable event-log and
        integration semantics are required.
\end{itemize}

OpenShift also supplies mature multi-tenancy, policy, secrets management,
operators, ingress, storage integration, observability, upgrades, ecosystem
compatibility, and operational practice. Charlotte must build corresponding
operational maturity where its use cases require it. The research question is
which mechanisms become smaller or more direct when the OS already has the
desired trust and execution model.

Durga can generate Kubernetes manifests, container builds, and policy as well as
Charlotte artifacts. Charlotte's proposed advantage is a
direct target of native cluster and capability objects, avoiding the
intermediate machine-like Linux userspace.

\section{Developer experience and the compatibility boundary}

If successful, the ordinary developer would work with:

\begin{itemize}
    \item BPMN and domain-specific annotations;
    \item typed task inputs, outputs, errors, and business-rule interfaces;
    \item a host-side test harness for business behavior;
    \item generated local simulations or integration tests;
    \item process-level diagnostics mapped to the BPMN diagram; and
    \item a build command that selects a verified backend.
\end{itemize}

The generator and platform team would own:

\begin{itemize}
    \item the Charlotte ABI and generated wrappers;
    \item serialization, schema evolution, correlation, and retry protocols;
    \item capability-request generation and policy integration;
    \item signing, provenance, and process-bundle construction;
    \item supervision, observation, deployment, and upgrade integration; and
    \item tested connectors to durable logs and external systems.
\end{itemize}

This division does not eliminate the need to learn Durga. It substitutes a
domain and process programming model for an operating-system programming model.
That is useful only if regeneration is predictable, generated code is
inspectable, errors name model elements, and debugging can
cross the boundary between rule code and runtime without exposing every kernel
detail.

Backend portability also requires restraint in the task API. A business rule
that assumes a local path, arbitrary outbound socket, process environment, or
Java-specific runtime facility is not automatically portable. The developer
contract should make portable operations easy, platform-specific dependencies
explicit, and unsupported dependencies visible at build time.

\section{Current implementation and next experiment}

The infrastructure experiment is complete enough to generate minimal
Charlotte-native activities, build signed ELFs, admit a multi-component release
atomically, resolve replica placement, pull artifacts through RustFS-backed S3,
and observe exact-generation readiness. The next experiment adds one complete
typed business-rule path and one transactional Kafka step, then binds their
process semantics in a signed bundle.

\subsection{Experiment input}

Use a small BPMN process containing:

\begin{itemize}
    \item a start event;
    \item one service activity with a typed request and response;
    \item one successor or completion event; and
    \item an explicit annotation for any clock, log, or external service the
        activity requires.
\end{itemize}

The handwritten Rust module should implement one pure or nearly pure business
function. It must not call a Charlotte syscall or manually construct an IPC
message.

\subsection{Generated output}

The next \texttt{charlotte} increment should add:

\begin{enumerate}
    \item stable request and response schema identities with serialization code;
    \item a typed business-handler contract used unchanged by host-side unit
        tests and the generated Charlotte wrapper;
    \item process-instance, activity, correlation, and attempt identities in
        calls and observations;
    \item a transactional-step binding that invokes the handler without
        exposing its Kafka connector capability;
    \item retained provenance compatible with the existing CLS2 metadata; and
    \item a signed semantic bundle binding the BPMN digest, activity and schema
        identities, generated artifacts, communication edges, capability plan,
        and executable deployment descriptors.
\end{enumerate}

A compact fixture with one transactional activity and one successor is
sufficient to test this boundary. Its components may use fixed or replicated
placement over the available test cluster; guest count is incidental. Fusion,
native durable process recovery, and general process-version migration remain
later experiments.

\subsection{Execution and verification}

The experiment should verify the following end-to-end properties:

\begin{enumerate}
    \item Regenerating the backend does not modify the developer-owned business
        rule.

    \item The generated ELF is signed, name-bound, digest-pinned, loaded in a
        fresh Charlotte address space, and invocable through the distributed
        name service.

    \item The activity receives only the generated capabilities. An attempted
        undeclared communication fails because no authority is present, not
        because a network deny rule happens to match.

    \item Typed input reaches the rule. Its typed result reaches the generated
        successor or test driver.

    \item Activity-entered, completed, and failed observations retain the BPMN
        process and activity identities.

    \item The existing stateless assignment handoff can move the generated
        artifact between nodes without changing its logical activity identity.

    \item The process-bundle metadata can be independently checked against the
        BPMN, schemas, generated source, executable digest, and retained
        provenance.
\end{enumerate}

The experiment should also record the amount and kind of Charlotte-specific
code that remains handwritten. Its success criterion is a developer-owned
business rule with no Charlotte programming, surrounded by a generated and
trusted runtime layer small enough to inspect and verify.

\subsection{Questions the experiment should answer}

The experiment is intended to expose design facts early:

\begin{itemize}
    \item What is the minimum Charlotte ABI needed by a generated activity?
    \item Which capability requirements can be inferred, and which require an
        explicit annotation or policy decision?
    \item Which parts of Durga's current semantic contract accidentally depend
        on Kafka?
    \item Which schema and correlation identities must cross the wrapper
        boundary?
    \item How should a business failure differ from a wrapper, transport, or
        node failure?
    \item Can host-side tests and Charlotte execution use the same rule code
        without conditional platform logic?
    \item Which identities are required to map runtime telemetry back to the
        BPMN model and signed build?
\end{itemize}

Answers to these questions determine whether a larger backend is justified.

\section{Remaining work}

Substantial work remains in both projects:

\begin{itemize}
    \item A substrate-neutral Durga semantic intermediate representation must
        distinguish process semantics from the present Kafka realization.

    \item A stable generated-task contract, schema evolution rules, and explicit
        delivery and failure semantics are required before backends can be
        compared meaningfully.

    \item The generated userspace ABI now covers launch, capability pickup, and
        lifecycle integration. Typed activity invocation, schema evolution,
        process observation, and durable recovery still need stable contracts.

    \item The Charlotte backend generates reviewable S3/Kafka profile requests
        without embedding credentials in developer-owned business code. Secret
        delivery, trust-anchor rotation, and profile replacement still need a
        production provisioning path.

    \item Generated capability plans and signed descriptor grants establish the
        request and admission shape. Inference rules, organizational review,
        policy evolution, and audit require a production model. The generator
        requests authority; a separate policy owner grants it.

    \item Process-bundle format, signatures, artifact and schema identity,
        provenance, and reproducible generation remain to be designed.

    \item The Charlotte cluster now has shared central-S3 artifact fetching and
        replica-set rollout observation and DSR routing. It still needs
        production administration authority, resource-aware placement,
        cluster-wide telemetry, and rolling update/rollback control
        (\autoref{chap:cluster-vision}).

    \item Native Charlotte messaging needs defined delivery, back-pressure,
        cancellation, deduplication, and recovery semantics before it can
        replace Kafka on any process edge that depends on them.

    \item Durable state, process-instance recovery, timers, human tasks,
        external side effects, and rolling process-version migration remain
        application-runtime problems even on a cluster OS.

    \item Debugging, local simulation, test fixtures, incident analysis, and
        operational explanations must be usable without requiring application
        teams to understand Charlotte internals.
\end{itemize}

The integration makes these questions testable one layer at a time while
keeping application semantics distinct from substrate mechanics.

\section{Next experiment}

The next falsifiable step is to take the generated prototype through a
multi-activity release. The experiment should determine whether its developer
can remain concerned with the business rule while generated code and the
cluster account for deployment, authority, and the operating environment.

\endinput
